% !TEX root = ../../../main.tex
% Block card: preprocessing.detrend -- "Baseline removal" (Part I, the blocks chapter).
\cardfield{What it does}
Takes away the slowly wandering level of a signal by subtracting a moving average (or a single
straight line), so that what is left sits around zero.

\cardfield{When to use it}
Before almost any detector or symbolic encoding on a channel that drifts. It is like measuring
the height of waves from the water's current level rather than from the sea floor: the moving
average stands in for the current level. It is the step the drop-detection template starts
with, in its \emph{rolling\_mean\_nearest} form. If large drops or spikes dominate the
channel, \blockref{preprocessing.beads} disturbs them less.

\cardfield{What goes in, what comes out}
In goes a span of one channel. Out comes a signal of the same length and rate: the original
minus its local average. In the \emph{rolling\_z} form it is also divided by the local spread,
so it comes out without units.

\cardfield{Settings that matter}
\emph{mode} chooses the method: \emph{rolling\_mean} (the default) subtracts a moving average;
\emph{rolling\_mean\_nearest} is the drop detector's own version, which differs only near the
ends of the span; \emph{rolling\_z} also rescales to a local z-score; \emph{linear} subtracts
one straight line fitted to the whole span. \emph{window\_s} is how many seconds the moving
average covers (ten minutes by default); it should be much longer than the events you want to
keep.

\cardfield{Watch out for}
A big drop drags the moving average down with it, so after subtraction the trace overshoots
upwards on either side of the drop, and an event longer than the window is partly removed.
Missing samples stay missing in the rolling modes, but \emph{linear} fails on them.
